\documentclass[10pt]{article}
\usepackage[T1]{fontenc}
\usepackage[utf8]{inputenc}
\usepackage[french]{babel}
\usepackage{amsmath,amssymb,mathrsfs,mathpazo}
\usepackage{geometry}
\geometry{a4paper,margin=2cm}
\usepackage{graphicx,array,tabularx,tabularray,colortbl,multirow,multicol,booktabs,enumitem,pifont,ulem,fancybox,diagbox,cancel}
\usepackage{pstricks,pstricks-add,pst-plot,pst-tree,pst-node,pst-text,pst-func,pst-math,pst-bspline,pst-3dplot,pst-eucl,pst-solides3d}
\let\pscalculate\undefined
\usepackage{pst-bezier}
\usepackage{tikz,pgfplots,tkz-tab}
\usetikzlibrary{arrows,arrows.meta,patterns,calc,positioning,trees,decorations.markings,shapes.geometric,shapes.misc,shapes.arrows,intersections,decorations.pathmorphing,decorations.pathreplacing,backgrounds,shadows,babel}
\pgfplotsset{compat=1.18}
\usepgfplotslibrary{fillbetween}
\usepackage[np]{numprint}
\usepackage{listings,makecell,pas-tableur,twoopt}
\let\barre\undefined
\usepackage[pstricks]{tablvar}
\definecolor{darkgreen}{RGB}{0,100,0}
\providecommand{\red}{\color{red}}
\providecommand{\blue}{\color{blue}}
\providecommand{\green}{\color{green}}
\providecommand{\black}{\color{black}}
\providecommand{\white}{\color{white}}
\providecommand{\gray}{\color{gray}}
\providecommand{\cyan}{\color{cyan}}
\providecommand{\magenta}{\color{magenta}}
\providecommand{\yellow}{\color{yellow}}
\providecommand{\violet}{\color{violet}}
\providecommand{\orange}{\color{orange}}
\providecommand{\euro}{\texteuro}
\providecommand{\EUR}{\texteuro}
\providecommand{\ieme}{\textsuperscript{e}}
\providecommand{\degres}{\ensuremath{^\circ}}
\providecommand{\degree}{\ensuremath{^\circ}}
\providecommand{\vv}[1]{\overrightarrow{#1}}
\providecommand{\Coord}[1]{#1}
\providecommand{\decosix}{}
\providecommand{\starredbullet}{\ensuremath{\bullet}}
\pagestyle{empty}
\setlength{\parindent}{0pt}
\usepackage{adjustbox}\DeclareUnicodeCharacter{2236}{:}\DeclareUnicodeCharacter{2190}{\ensuremath{\leftarrow}}
\providecommand{\R}{\mathbb{R}}
\providecommand{\N}{\mathbb{N}}
\providecommand{\Z}{\mathbb{Z}}
\providecommand{\Q}{\mathbb{Q}}
\providecommand{\D}{\mathbb{D}}
\providecommand{\C}{\mathbb{C}}
\providecommand{\e}{{\mathrm{e}}}
\providecommand{\vect}[1]{\overrightarrow{#1}}
\providecommand{\vectt}[1]{\overrightarrow{\mathrm{#1}}}
\providecommand{\vvt}[1]{\overrightarrow{\mathrm{#1}}}
\providecommand{\np}[2][]{{#2~#1}}
\providecommand{\vv}[1]{{\overrightarrow{#1}}}
\renewcommand{\d}{{\mathrm{d}}}
\providecommand{\degres}{{^\circ}}
\providecommand{\degre}{{^\circ}}
\providecommand{\blue}{\color{blue}}
\providecommand{\red}{\color{red}}
\providecommand{\enskip}{\quad}
\providecommand{\llbracket}{\lbrack\!\lbrack}
\providecommand{\rrbracket}{\rbrack\!\rbrack}
\providecommand{\euro}{{\text{€}}}
\providecommand{\no}{{\text{n°}}}
\providecommand{\interval}[2]{{[#1;#2]}}
\providecommand{\intervFO}[2]{{[#1;#2[}}
\providecommand{\dotfill}{\ldots}

\providecommand{\ds}{\displaystyle}
\providecommand{\textding}[1]{#1}
\providecommand{\Oij}{\(\mathrm{O};\vec{\imath},\vec{\jmath}\)}
\providecommand{\Oijk}{\(\mathrm{O};\vec{\imath},\vec{\jmath},\vec{k}\)}
\providecommand{\Ouv}{\(\mathrm{O};\vec{u},\vec{v}\)}
\providecommand{\intervFF}[2]{\ensuremath{[#1;#2]}}
\renewcommand{\euro}{\ifmmode\text{\texteuro}\else\texteuro\fi}
\DeclareUnicodeCharacter{00B0}{\ifmmode^{\circ}\else\textdegree\fi}
\DeclareUnicodeCharacter{20AC}{\ifmmode\text{\texteuro}\else\texteuro\fi}
% Énoncés : Éducation nationale. Transcription et corrigés : APMEP, Denis Vergès et contributeurs. Découpage et classement : CoopMaths, Sébastien Lozano, Éric Elter et contributeurs. Adaptation de la présentation : Maths Explorer.
% CC BY-SA : https://coopmaths.fr/www/about/
% https://coopmaths.fr/alea/static/bac/2025/tex/bac_2025_11_sujet2_ameriquesud_2.tex
\begin{document}
\noindent\textit{APMEP / CoopMaths. CC BY-SA. Présentation adaptée par Maths Explorer.}\par\medskip
\medskip

Cet exercice est un questionnaire à choix multiple. Pour chaque question, une seule des trois propositions est exacte.

Indiquer sur la copie le numéro de la question et la lettre de la proposition
choisie.

Aucune justification n'est demandée.

\medskip

\emph{Pour chaque question, une réponse exacte rapporte un point.\\
Une réponse fausse, une réponse multiple ou l'absence de réponse ne rapporte ni n'enlève de point.}

\medskip

Dans toutes les questions suivantes, l'espace est rapporté à un repère orthonormé.

\medskip

\begin{enumerate}
\item On considère la droite $\Delta_1$ de représentation paramétrique 
$\left\{\begin{array}{l c l}
x &=& 1 - 3t\\
y &=& 4 + 2t\\
z &=&\phantom{4 +2}t
\end{array}\right.$, où $t \in  \R$
ainsi que la droite $\Delta_2$ de représentation paramétrique 
$\left\{\begin{array}{l c l}
x& =& -4 + \phantom{2}s\\
y& =& \phantom{-}2 + 2s\\
z& =& - 1 + \phantom{2}s
\end{array}\right.$, où $s \in \R$.
\begin{enumerate}
\item Les droites $\Delta_1$ et $\Delta_2$ sont parallèles.
\item Les droites $\Delta_1$ et $\Delta_2$ sont orthogonales.
\item Les droites $\Delta_1$ et $\Delta_2$ sont sécantes.
\end{enumerate}
\item On considère la droite $d$ de représentation paramétrique $\left\{\begin{array}{l c l}x&=&1 + \phantom{2}t\\
y &=& 3 - \phantom{2}t\\
z &=& 1 + 2t
\end{array}\right.$, où $t \in \R$,
et le plan $P$ d'équation cartésienne : $4x + 2y - z + 3 = 0.$
	\begin{enumerate}
		\item La droite $d$ est incluse dans le plan $P$.
		\item La droite $d$ est parallèle strictement au plan $P$.
		\item La droite $d$ est sécante au plan $P$.
	\end{enumerate}
\item On considère les points A(3~;~2~;~1), B(7~;~3 ~;~1), C$(-1~;~4~;~5)$ et D$(- 3~;~3~;~5)$.
	\begin{enumerate}
		\item Les points A, B, C et D ne sont pas coplanaires.
		\item Les points A, B et C sont alignés.
		\item $\vect{\text{AB}}$ et $\vect{\text{CD}}$ sont colinéaires.
	\end{enumerate}
\item On considère les plans $Q$ et $Q'$ d'équation cartésienne respective $3x - 2y + z + 1 = 0$ et $4x + y - z + 3 = 0$.
	\begin{enumerate}
		\item Le point R$(1~;~1~;~-2)$ appartient aux deux plans.
		\item Les deux plans sont orthogonaux.
		\item Les deux plans sont sécants avec pour intersection la droite de représentation
paramétrique $\left\{\begin{array}{l c l}
x&=& \phantom{17}t\\
y& =& \phantom{1}7t + 4 \\
z& =& 11t + 7
\end{array}\right.$, où $t \in \R$.
		\end{enumerate}
\end{enumerate}

\vspace{0,5cm}

\end{document}
